\clearpage
\appendix

\section{Kolmogorov-Smirnov Test Details}
\label{sec:app-ks-test}

In our setting, let $F_U(x)$ comprising $n$ samples and $F_R(x)$ comprising $m$ samples be the empirical CDF of the unlearned and retain models, respectively.
Then, the KS-Test computes a statistic $D_{n,m} = \sup_{x} \vert F_U(x) - F_R(x) \vert$.

The null hypothesis, stating that the two sets of samples are drawn from the same distribution, is rejected at a chosen significance level $\alpha$ if the following inequality holds.
\begin{equation}
    D_{n,m} > c(\alpha) \sqrt{\frac{n + m}{nm}},
\end{equation}
where $c(\alpha)$ is the critical value of that significance level.
\begin{equation}
    c(\alpha) = \sqrt{-\ln\left(\frac{\alpha}{2}\right) \cdot \frac12}.
\end{equation}

The $p$-value is then defined as the minimal alpha for which the inequality holds, or the smallest value at which we can reject the null hypotheses.
Forget quality is hence, a measure of the confidence that the distributions of Truth Ration values over the forget set from two models are the same.

\section{Hyperparameters}
\label{sec:app-hyperparams}
We trained using AdamW with weight decay of $0.01$ and $0$. The learning rate is fixed to be 1e-5. For both finetuning and unlearning, we fix the epoch number to be 5, and we incorporate a linear warmup in the first epoch. 
We experiment with various learning rates including $10^{-5}$, $10^{-6}$, and $5\cdot10^{-7}$ and find that $10^{-5}$ is a good choice for the baseline methods in our experiments. 
In particular, $10^{-6}$ is too small of a learning rate to see appreciable updates to the model weights.

\section{Preference Strings}
\label{sec:app-idk-strings}

\input{tables/idk-strings}

\section{Continued Unlearning}
\label{sec:app-continued-unlearning}

In the main experiments of this paper, we limit unlearning to five epochs, but one might wonder how things progress given more time to unlearn.
We test forgetting $1\%$ of the data with Phi-1.5 and show that continued unlearning does not help with these baseline methods, see Figure~\ref{fig:12epochs}.

\begin{figure}[H]
    \centering
    \includegraphics[width=0.5\textwidth]{figures/scatterplots/wd0.01-forget01-12epochs.pdf}
    \caption{Zoomed in plots of extended unlearning trajectories (10 epochs) on the $1\%$ forget set.}
    \label{fig:12epochs}
\end{figure}

\section{Sanity Checks}
\label{sec:app-sanity}

\begin{table}[t!]
\centering
\caption{Model comparisons using KS-Test $p$-values for Llama-2-7B Models (WD $= 0.00$). We compare retain models finetuned with $90\%$, $95\%$, and $99\%$ of the data. We test the Truth Ratio distributions over both retain data and forget data. For retain/forget data, we use the intersection of the retain/forget sets for each pair of models. All of these $p$-values are high indicating that the KS-Test accurately captures the similarity we know these models have over each of these datasets.}
\label{tab:llama-pvals-1}
\footnotesize
\begin{tabular}{lcccccccccccc}
\toprule
                    &    & Retain 90 & Retain 95 & Retain 99   \\
\midrule
                    &   Retain 90 & 1             & 0.9414        & 0.8483        \\
Retain Data         &   Retain 95 & -             & 1             & 0.9705        \\
                    &   Retain 99 & -             & -             & 1             \\
\midrule
                    &   Retain 90 & 1             & 0.8655        & 0.7659       \\
Forget Data         &   Retain 95 & -             & 1             & 0.9900       \\
                    &   Retain 99 & -             & -             & 1            \\
\bottomrule
\end{tabular}
\end{table}

\begin{table}[t!]
\centering
\caption{Model comparisons using KS-Test $p$-values for Llama-2-7B Models (WD $= 0.00$). We compare retain models finetuned with $90\%$, $95\%$, and $99\%$ of the data to a model finetuned on all the \tofu{} data, a pretrained base model, and a random model. We test the Truth Ratio distributions over both retain data and forget data. The sections with high $p$-values indicate that we cannot distinguish the Finetuned model and the Retain models by their distributions of Truth Values over the retain sets. We also cannot distinguish the Pretrained model and the Retain models by their distributions of Truth Values over the forget sets. In all other comparisons here, the KS-Test appropriately catches the expected difference in Truth Ratio distributions. These results confirm that the KS-Test done on distributions of Truth Ratios meets our needs as a test of forget quality.}
\label{tab:llama-pvals-2}
\footnotesize
\begin{tabular}{lcccccccccccc}
\toprule
         & & Finetuned & Pretrained & Random \\
\midrule
                & Retain 90 & 0.9705      & 9.21E-31      & 2.42E-66  \\
Retain Data     & Retain 95 & 0.9879      & 1.41E-32      & 2.94E-69  \\
                & Retain 99 & 0.9003      & 4.07E-32      & 2.94E-69  \\ 
\midrule
                & Retain 90  & 1.10E-19     & 0.0031       & 2.43E-19  \\
Forget Data     & Retain 95  & 4.73E-15     & 0.0297       & 2.96E-13  \\
                & Retain 99  & 5.04E-04     & 0.1650       & 5.04E-04  \\ 
\bottomrule
\end{tabular}
\end{table}






We verify that our metrics for Model Utility and Forget Quality have some desirable properties. 
In Tables~\ref{tab:llama-pvals-1} and \ref{tab:llama-pvals-2}, we show the $p$-values for the KS-Tests that confirm all of our expectations enumerated below and validate this choice of metric.
These tables have figures from Llama-2-7B tests, but the same trends hold for Phi-1.5.

First, Model Utility should meet the following natural expectations.
\begin{enumerate}
    \item Model Utility should be high for a pretrained model (one that has never been finetuned on \tofu{} data).
    \item Model Utility should be low for a model with random weights.
\end{enumerate}

Additionally, Forget Quality is measured using a statistical test on Truth Ratio values, and so we hope that this test meets the following expectations.
\begin{enumerate}
    \item The KS-Test performed on distributions of Truth Ratio values over the intersection of the three forget sets (from the 90-10, 95-5, and 99-1 splits) should produce high $p$-values when comparing any two retain models.
    \item The KS-Test performed on distributions of Truth Ratio values over the intersection of the three retain sets (from the 90-10, 95-5, and 99-1 splits) should produce high $p$-values when comparing any two retain models.
    \item The KS-Test performed on distributions of Truth Ratio values over the forget set should produce high $p$-values when comparing any retain model to a random model.
    \item The KS-Test performed on distributions of Truth Ratio values over the retain set should produce low $p$-values when comparing any retain model to a random model.
    \item The KS-Test performed on distributions of Truth Ratio values over the forget set should produce high $p$-values when comparing any retain model to a pretrained model.
    \item The KS-Test performed on distributions of Truth Ratio values over the retain set should produce low $p$-values when comparing any retain model to a pretrained model.
    \item The KS-Test performed on distributions of Truth Ratio values over the forget set should produce low $p$-values when comparing any retain model to a finetuned model (finetuned on all the \tofu{} data and without any unlearning).
    \item The KS-Test performed on distributions of Truth Ratio values over the retain set should produce high $p$-values when comparing any retain model to a finetuned model (finetuned on all the \tofu{} data and without any unlearning).
\end{enumerate}

\section{Knowledge Entanglement}
\label{sec:app-knowledge-entanglement}

One of the challenges of unlearning comes from knowledge entanglement---when we try to make a model forget about one thing, it also tends to forget other things unexpectedly. 
This phenomenon is similar to catastrophic forgetting in continual learning \citep{mccloskey1989catastrophic}. 
In Figures~\ref{fig:knowledge-entanglement-0}-\ref{fig:knowledge-entanglement-23}, we show this phenomenon in different models and unlearning algorithms. 
In \Cref{fig:rouge_oracle_llama}, even with access to the oracle model or retain set, model generation on all four sets still has a decreasing ROUGE, especially the dataset that relate to authors. This suggests the existence of knowledge entanglement, showing why unlearning is hard.
Consider the case of unlearning the 5\% forget set with Gradient Difference on  Llama-2-7B, Figure~\ref{fig:knowledge-entanglement-7}.
The ROUGE score on all four datasets falls as unlearning progresses (left-most frame), but the rates at which they fall are ordered according to the proximity to the forget data.
(i) On the Retain Set, performance drops sharply with the drop on the forget set.
(ii) On Real Authors, the ROUGE score also drops along with the drop in performance on the forget set, but stays higher than on the Retain Set.
(iii) Finally, performance on World Facts stays relatively unchanged.

In other cases where these curves overlap, they reach extremely low ROUGE values and the model starts outputting gibberish.
This suggests the existence of \emph{knowledge
entanglement}, supporting that our choice of having multiple evaluation datasets is important for a holistic assessment of unlearning.

\twocolumn

\begin{figure}[tb]
    \centering
    \includegraphics[width=\columnwidth]{figures/all_metrics/llama/1GPU_grad_ascent_1e-05_forget01_all3metric.pdf}
    \caption{Unlearn Llama-2-7B with gradient ascent on $1\%$ forget set.}
    \label{fig:knowledge-entanglement-0}
\end{figure}
\begin{figure}[tb]
    \centering
    \includegraphics[width=\columnwidth]{figures/all_metrics/llama/1GPU_grad_ascent_1e-05_forget05_all3metric.pdf}
    \caption{Unlearn Llama-2-7B with gradient ascent on $5\%$ forget set.}
    \label{fig:knowledge-entanglement-1}
\end{figure}
\begin{figure}[tb]
    \centering
    \includegraphics[width=\columnwidth]{figures/all_metrics/llama/1GPU_grad_ascent_1e-05_forget10_all3metric.pdf}
    \caption{Unlearn Llama-2-7B with gradient ascent on $10\%$ forget set.}
    \label{fig:knowledge-entanglement-2}
\end{figure}

\begin{figure}[tb]
    \centering
    \includegraphics[width=\columnwidth]{figures/all_metrics/llama/1GPU_dpo_1e-05_forget01_all3metric.pdf}
    \caption{Unlearn Llama-2-7B with preference optimization on $1\%$ forget set.}
    \label{fig:knowledge-entanglement-3}
\end{figure}
\begin{figure}[tb]
    \centering
    \includegraphics[width=\columnwidth]{figures/all_metrics/llama/1GPU_dpo_1e-05_forget05_all3metric.pdf}
    \caption{Unlearn Llama-2-7B with preference optimization on $5\%$ forget set.}
    \label{fig:knowledge-entanglement-4}
\end{figure}
\begin{figure}[tb]
    \centering
    \includegraphics[width=\columnwidth]{figures/all_metrics/llama/1GPU_dpo_1e-05_forget10_all3metric.pdf}
    \caption{Unlearn Llama-2-7B with preference optimization on $10\%$ forget set.}
    \label{fig:knowledge-entanglement-5}
\end{figure}

\begin{figure}[tb]
    \centering
    \includegraphics[width=\columnwidth]{figures/all_metrics/llama/1GPU_KL_1e-05_forget01_all3metric.pdf}
    \caption{Unlearn Llama-2-7B with gradient difference on $1\%$ forget set.}
    \label{fig:knowledge-entanglement-6}
\end{figure}
\begin{figure}[tb]
    \centering
    \includegraphics[width=\columnwidth]{figures/all_metrics/llama/1GPU_KL_1e-05_forget05_all3metric.pdf}
    \caption{Unlearn Llama-2-7B with gradient difference on $5\%$ forget set.}
    \label{fig:knowledge-entanglement-7}
\end{figure}
\begin{figure}[tb]
    \centering
    \includegraphics[width=\columnwidth]{figures/all_metrics/llama/1GPU_KL_1e-05_forget10_all3metric.pdf}
    \caption{Unlearn Llama-2-7B with gradient difference on $10\%$ forget set.}
    \label{fig:knowledge-entanglement-8}
\end{figure}

\begin{figure}[tb]
    \centering
    \includegraphics[width=\columnwidth]{figures/all_metrics/llama/1GPU_oracle_1e-05_forget01_all3metric.pdf}
    \caption{Unlearn Llama-2-7B with KL Minimization on $1\%$ forget set.}
    \label{fig:knowledge-entanglement-9}
\end{figure}
\begin{figure}[tb]
    \centering
    \includegraphics[width=\columnwidth]{figures/all_metrics/llama/1GPU_oracle_1e-05_forget05_all3metric.pdf}
    \caption{Unlearn Llama-2-7B with KL Minimization on $5\%$ forget set.}
    \label{fig:knowledge-entanglement-10}
\end{figure}
\begin{figure}[tb]
    \centering
    \includegraphics[width=\columnwidth]{figures/all_metrics/llama/1GPU_oracle_1e-05_forget10_all3metric.pdf}
    \caption{Unlearn Llama-2-7B with KL Minimization on $10\%$ forget set.}
    \label{fig:knowledge-entanglement-11}
\end{figure}


\begin{figure}[tb]
    \centering
    \includegraphics[width=\columnwidth]{figures/all_metrics/phi/1GPU_grad_ascent_1e-05_forget01_all3metric.pdf}
    \caption{Unlearn Phi with gradient ascent on $1\%$ forget set.}
    \label{fig:knowledge-entanglement-12}
\end{figure}
\begin{figure}[tb]
    \centering
    \includegraphics[width=\columnwidth]{figures/all_metrics/phi/1GPU_grad_ascent_1e-05_forget05_all3metric.pdf}
    \caption{Unlearn Phi with gradient ascent on $5\%$ forget set.}
    \label{fig:knowledge-entanglement-13}
\end{figure}
\begin{figure}[tb]
    \centering
    \includegraphics[width=\columnwidth]{figures/all_metrics/phi/1GPU_grad_ascent_1e-05_forget10_all3metric.pdf}
    \caption{Unlearn Phi with gradient ascent on $10\%$ forget set.}
    \label{fig:knowledge-entanglement-14}
\end{figure}

\begin{figure}[tb]
    \centering
    \includegraphics[width=\columnwidth]{figures/all_metrics/phi/1GPU_dpo_1e-05_forget01_all3metric.pdf}
    \caption{Unlearn Phi with preference optimization on $1\%$ forget set.}
    \label{fig:knowledge-entanglement-15}
\end{figure}
\begin{figure}[tb]
    \centering
    \includegraphics[width=\columnwidth]{figures/all_metrics/phi/1GPU_dpo_1e-05_forget05_all3metric.pdf}
    \caption{Unlearn Phi with preference optimization on $5\%$ forget set.}
    \label{fig:knowledge-entanglement-16}
\end{figure}
\begin{figure}[tb]
    \centering
    \includegraphics[width=\columnwidth]{figures/all_metrics/phi/1GPU_dpo_1e-05_forget10_all3metric.pdf}
    \caption{Unlearn Phi with preference optimization on $10\%$ forget set.}
    \label{fig:knowledge-entanglement-17}
\end{figure}

\begin{figure}[tb]
    \centering
    \includegraphics[width=\columnwidth]{figures/all_metrics/phi/1GPU_KL_1e-05_forget01_all3metric.pdf}
    \caption{Unlearn Phi with gradient difference on $1\%$ forget set.}
    \label{fig:knowledge-entanglement-18}
\end{figure}
\begin{figure}[tb]
    \centering
    \includegraphics[width=\columnwidth]{figures/all_metrics/phi/1GPU_KL_1e-05_forget05_all3metric.pdf}
    \caption{Unlearn Phi with gradient difference on $5\%$ forget set.}
    \label{fig:knowledge-entanglement-19}
\end{figure}
\begin{figure}[tb]
    \centering
    \includegraphics[width=\columnwidth]{figures/all_metrics/phi/1GPU_KL_1e-05_forget10_all3metric.pdf}
    \caption{Unlearn Phi with gradient difference on $10\%$ forget set.}
    \label{fig:knowledge-entanglement-20}
\end{figure}

\begin{figure}[tb]
    \centering
    \includegraphics[width=\columnwidth]{figures/all_metrics/phi/1GPU_oracle_1e-05_forget01_all3metric.pdf}
    \caption{Unlearn Phi with KL Minimization on $1\%$ forget set.}
    \label{fig:knowledge-entanglement-21}
\end{figure}
\begin{figure}[tb]
    \centering
    \includegraphics[width=\columnwidth]{figures/all_metrics/phi/1GPU_oracle_1e-05_forget05_all3metric.pdf}
    \caption{Unlearn Phi with KL Minimization on $5\%$ forget set.}
    \label{fig:knowledge-entanglement-22}
\end{figure}
\begin{figure}[tb]
    \centering
    \includegraphics[width=\columnwidth]{figures/all_metrics/phi/1GPU_oracle_1e-05_forget10_all3metric.pdf}
    \caption{Unlearn Phi with KL Minimization on $10\%$ forget set.}
    \label{fig:knowledge-entanglement-23}
\end{figure}
